\bookchapter{Capítulo 38: As cinco grandes regras de pensamento da programação com IA}

\chapterintro{Se deste livro você só puder guardar um capítulo, que seja este. As cinco regras são o coração da obra; tudo o que vem depois — contexto, requisitos, prompt, qualidade — é a aplicação delas em cada cenário. Fique tranquilo: nenhuma é difícil, mas todas merecem ser “internalizadas” com alguns projetos.}

A história começa de um jeito curioso: ao abrir o projeto, Andrej Karpathy anexou um CLAUDE.md de só 70 linhas com quatro princípios de colaboração com IA. O arquivo recebeu mais de 70.000 Stars no GitHub e virou a referência mais citada entre quem programa com IA. Não são slogans vazios: são hábitos mentais práticos, destilados da prática. Este livro acrescenta o quinto — iteração e verificação — completando as cinco regras.

\begin{tcbinfo}{Por que 70 linhas valem tanto?}
70.000 Stars? A maioria dos projetos no GitHub nunca chega perto disso. Que um arquivo de 70 linhas alcance esse número prova uma coisa: na era da IA, o mais raro não é código — é a capacidade de “pensar com clareza”. Essas 70 linhas são a versão concentrada desse pensamento.

\end{tcbinfo}

\booksection{Regra 1: Think Before Coding (pense antes de codar)}

É a regra mais contra-intuitiva das cinco. Se a IA escreve o código, por que “pensar antes”? Resposta simples: a qualidade da saída é proporcional à da entrada. Se você mesmo não sabe o que quer, a IA só te leva a um ciclo infinito de “muda e muda de novo” — o famoso “a IA se perde junto com você”.

“Pensar antes” significa responder três perguntas antes de dar a instrução; juntas, formam o “triângulo do pensamento”:

\begin{enumerate}
  \item Qual problema quero resolver? (objetivo)
  \item Como deve ser o resultado esperado? (critério de aceite)
  \item Quais restrições preciso respeitar? (limites)
\end{enumerate}

\begin{center}
\begin{tabularx}{\linewidth}{|X|X|X|}
\hline
\textbf{Dimensão} & \textbf{Entrada vaga (ruim)} & \textbf{Entrada clara (boa)} \tabularnewline
\hline
Definição do problema & Faz um site pra mim & Faz uma página de ferramentas pessoais com cartões de entrada para 3 miniapps de IA \tabularnewline
Resultado esperado & Que fique bonito & Layout responsivo, tema escuro, animação ao passar o mouse, funciona no celular \tabularnewline
Restrições & (não dito) & Sem frameworks, HTML+CSS+JS puro, deploy em arquivo único, sem backend \tabularnewline
\hline
\end{tabularx}
\end{center}

\begin{center}\small\color{gray}Comparação entre entrada vaga e clara para o mesmo pedido\end{center}

\begin{tcbwarn}{Atenção: pensar antes $\neq$ escrever um documento perfeito}
“Pensar antes” não é “fazer um PRD perfeito e só então começar”. É um processo rápido e organizado: 30 segundos para clarear as três perguntas valem mais do que 30 minutos de vaivém no chat. Em tarefas pequenas, as respostas podem ser curtas; em projetos grandes, viram um documento de requisitos completo (Capítulo 40).

\end{tcbwarn}

\booksection{Regra 2: Simplicity First (simplicidade em primeiro lugar)}

A maior armadilha da programação com IA não é “não conseguir fazer”, e sim “fazer demais”. Se você pede “um sistema de usuários”, a IA pode gerar autenticação, permissões, papéis, logs de auditoria, exportação… quando você só queria uma tela de login. Excesso de geração é a principal causa de código inchado.

“Simplicidade em primeiro lugar” é aplicar o princípio do “mínimo viável” a cada instrução: menos código para a função central, a solução mais simples para o problema, o mínimo de recursos para validar a ideia. Karpathy tem uma comparação precisa: “para calcular um preço com desconto, você não precisa de um motor de e-commerce.”

Na prática, significa alguns hábitos:

\begin{itemize}
  \item Peça à IA uma função de cada vez, não o sistema inteiro de uma vez;
  \item Prefira a biblioteca padrão e soluções simples a dependências complexas;
  \item Se cabe em 50 linhas, não deixe a IA escrever 500;
  \item Se a entrega passou do escopo, mande cortar o excesso sem medo.
\end{itemize}

\begin{center}
\begin{tabularx}{\linewidth}{|X|X|X|}
\hline
\textbf{Cenário} & \textbf{Solução exagerada (evite)} & \textbf{Solução simples (recomendada)} \tabularnewline
\hline
Armazenar dados & Montar MySQL + ORM & Começar com arquivo JSON ou SQLite \tabularnewline
Login de usuário & OAuth2.0 + JWT + papéis e permissões & Começar com validação simples de senha \tabularnewline
Framework de front-end & React + Redux + Router & Testar HTML + JavaScript puro \tabularnewline
Publicar online & Docker + CI/CD + K8s & Fazer rodar local primeiro \tabularnewline
\hline
\end{tabularx}
\end{center}

\begin{center}\small\color{gray}Soluções exageradas vs. simples em cenários comuns\end{center}

\booksection{Regra 3: Surgical Changes (modificações cirúrgicas)}

Ao iterar um projeto existente, cada mudança deve ser cirúrgica: mexa só onde precisa. Em programação com IA, isso é crucial — a IA adora “aproveitar e melhorar mais um pouco”: você pede para corrigir um bug e ela refaz metade do arquivo.

O coração da regra: toda instrução deve fixar o escopo e os limites da mudança. Diga “modifique só a função calculate\_discount”, “não mexa no resto”, “mantenha o estilo atual”. Se a IA ultrapassar o pedido, mande reverter o excesso na hora.

Karpathy ilustra com um exemplo de diff: pedir à IA para adicionar “preço de membro” à função de desconto do carrinho. A boa entrega altera só essa função e acrescenta poucas linhas; a má entrega refaz o arquivo inteiro — o código novo pode até parecer “mais elegante”, mas você não sabe se introduziu bugs. Em programação com IA, “parecer melhor” nunca é motivo para sair do pedido.

Instrução ruim (escopo vago)

\begin{lstlisting}
Otimiza a lógica de desconto deste projeto
\end{lstlisting}

Instrução boa (escopo preciso)

\begin{lstlisting}
Modifique apenas a função calculate\_discount em cart.py: | 1. Adicione lógica de preço de membro sobre o desconto atual | 2. Membros ganham 10\% extras | 3. Não altere a assinatura da função nem outras funções | 4. Mantenha o estilo atual de nomes de variáveis
\end{lstlisting}

\booksection{Regra 4: Goal-Driven Execution (execução orientada a objetivo)}

Esta regra responde a uma pergunta-chave: instruir a IA sobre “como fazer” ou sobre “o que fazer”? Karpathy prefere a segunda. Em vez de “crie uma classe, depois adicione métodos, depois escreva testes”, diga “preciso de uma ferramenta que leia um CSV e conte os valores únicos de cada coluna”. A IA escolhe o caminho; você valida o resultado.

Os benefícios são duplos. Para a IA, há mais autonomia para escolher a melhor implementação em vez de seguir seus passos intermediários; para você, a atenção sai do “como” e vai para o “está certo?” — exatamente o papel de um Conductor.

\begin{tcbwarn}{Atenção: orientado a objetivo $\neq$ deixar solto}
Depois da execução, você valida com rigor: o resultado atende ao pedido? Casos de borda foram tratados? Há overdesign óbvio? Se algo saiu errado, dê feedback concreto — não “refaça”, e sim “com lista vazia o código lança IndexError; adicione verificação de nulo”.

\end{tcbwarn}

\begin{center}
\begin{tabularx}{\linewidth}{|X|X|X|}
\hline
\textbf{Estilo de instrução} & \textbf{Exemplo} & \textbf{Característica do resultado} \tabularnewline
\hline
Imperativo (como fazer) & "Crie uma classe Python chamada DataParser, adicione um método parse que receba filename e use o módulo csv…" & A IA vira digitadora; seus passos intermediários podem não ser o caminho ótimo \tabularnewline
Por objetivo (o que fazer) & "Preciso de uma ferramenta: ler CSV, contar valores únicos por coluna e sair em JSON. Entregue a implementação completa." & A IA escolhe a melhor rota; você valida o resultado \tabularnewline
\hline
\end{tabularx}
\end{center}

\begin{center}\small\color{gray}Comparação entre instrução imperativa e orientada a objetivo\end{center}

\booksection{Regra 5: Iterate and Verify (iterar e verificar)}

É a quinta regra que este livro acrescenta às quatro de Karpathy. A crença central: a primeira entrega da IA quase nunca é a versão final. Programação eficaz com IA é um ciclo “gerar–verificar–corrigir”, não um único “pedir–receber”.

“Iterar e verificar” não é um “está errado, tenta de novo” passivo; é um fluxo estruturado. A cada entrega, valide em três camadas:

\begin{itemize}
  \item 1ª · Funcional: roda? Entradas e saídas estão corretas?
  \item 2ª · Borda: e com entrada vazia, gigante ou inválida?
  \item 3ª · Arquitetura: a estrutura faz sentido? Há hardcoded? Dá para evoluir?
\end{itemize}

Definir a meta

A IA gera código

Verificação funcional

Aprovado

Verificação de casos limite

Aprovado

Revisão da arquitetura

Aprovado — vai para a próxima etapa

Reprovado

Dar feedback específico

Regenerar

\begin{center}\small\color{gray}Figura 2: ciclo de iteração e verificação na programação com IA\end{center}

O ponto crucial é o “feedback específico”. Em vez de “está errado, refaça”, diga “com lista vazia a função lança IndexError; adicione verificação na linha 15”. Feedback preciso permite correção cirúrgica; feedback vago só faz a IA adivinhar de novo.

\booksection{Harness Engineering: Agente = Modelo + Harness}

Por volta de meados de 2026, o conceito Harness Engineering começou a popularizar-se na comunidade. A fórmula é simples:

A capacidade de um agent = capacidade do modelo × qualidade do harness

A capacidade do modelo é o teto intelectual da IA; o Harness é o “ambiente de trabalho” em volta do modelo — gestão de contexto, integração de ferramentas, orquestração de fluxo, sistema de memória etc.

Isso explica uma dúvida comum: por que, com o mesmo GPT-4 ou Claude, uns fazem projetos incríveis e outros saem com código medíocre? A diferença não está no modelo, e sim no Harness. Um bom Harness tem: especificação clara do projeto, ferramentas bem integradas, contexto organizado e fluxo de verificação sólido. O modelo define o teto; o Harness define quanto dele você alcança.

\begin{tcbinfo}{A metáfora do estagiário genial}
Imagine o modelo como um estagiário genial: altíssimo QI, mas sem documento de contexto, sem ferramentas, sem prazo e sem critérios de aceite, não entrega nada de bom. Harness Engineering é a engenharia de “como montar o ambiente de trabalho do estagiário genial”.

\end{tcbinfo}

Os próximos capítulos ensinam a construir esse Harness: Capítulo 39, contexto (memória do projeto); 40, requisitos (PRD); 43, prompt em camadas; 45, garantia de qualidade (verificação). Juntos, formam o seu Harness.

\begin{tcbpractice}{Mão na massa: use o "triângulo do pensamento" na sua próxima ideia}
Escolha um projetinho que você adia há tempo (uma página de contagem regressiva, um script de lembrete de beber água). Em 3 minutos responda: (1) que problema resolve? (2) como fica pronto? (3) quais restrições (tempo, tecnologia, ambiente)? Escreva as respostas em um parágrafo. Parabéns: você acabou de fazer o seu primeiro “Think Before Coding” — e esse parágrafo já é uma instrução de qualidade.

\end{tcbpractice}

\begin{tcbsummary}{Resumo do capítulo}
\begin{itemize}
  \item As cinco regras vêm do CLAUDE.md de 70 linhas de Karpathy (70 mil Stars); este livro acrescenta a quinta: "iterar e verificar"
  \item Think Before Coding: responda três perguntas — o que fazer? como fica o resultado? quais limites?
  \item Simplicity First: a maior armadilha é "fazer demais" — para calcular desconto não se precisa de um e-commerce
  \item Surgical Changes: mudanças cirúrgicas, sem o "aproveitar e melhorar" da IA
  \item Goal-Driven Execution: diga o que fazer, não como fazer — deixe a IA escolher o caminho
  \item Iterate and Verify: três camadas (funcional/borda/arquitetura) + feedback concreto, em ciclo
  \item Harness Engineering: capacidade = modelo × Harness; trate a IA como "estagiário genial" e monte o ambiente
\end{itemize}

\end{tcbsummary}
